\chapter{Performance: Where Code and Data Live}
\label{ch:perf}\index{performance}\index{memory map}\index{IRAM}\index{XIP}

On a chip with a uniform flat memory, performance is mostly about the algorithm. On
the ESP32-S3 it is also about \emph{placement}: the same loop runs at different
speeds depending on whether its instructions come from on-chip SRAM or are fetched
from external flash through a cache, and the same buffer is or isn't reachable by
DMA depending on which memory it lives in. This chapter collects the placement facts
scattered through the rest of the book into one picture, so you can reason about
speed deliberately instead of by surprise.

\section{The memory map, read for speed}\index{IROM}\index{DROM}

The linker scripts pin every region to a fixed address (the full table is in
Chapter~\ref{ch:bootloader}); read for performance, they fall into three speeds:

\begin{center}
\begin{tabular}{l l l}
\toprule
\textbf{Region} & \textbf{Base} & \textbf{Speed} \\ \midrule
IRAM (code in SRAM)   & \texttt{16\#4037\_8000\#} & full speed, always \\
DRAM (data in SRAM)   & \texttt{16\#3FC9\_0000\#} & full speed, always \\
IROM (code in flash)  & \texttt{16\#4200\_0000\#} & cache hit fast, miss = flash read \\
DROM (rodata in flash)& \texttt{16\#3C00\_0000\#} & cache hit fast, miss = flash read \\
PSRAM (external RAM)  & \texttt{16\#3D00\_0000\#} & cached, slower than SRAM, no DMA \\
\bottomrule
\end{tabular}
\end{center}

The bulk of a program (all its \texttt{.text} and read-only data) lives in
flash and executes \emph{in place} (XIP): the CPU fetches through the instruction
cache that \texttt{start.S} configures, and on a hit it runs at full speed. A cache
\emph{miss} pays a SPI-flash read of the missing line. For straight-line and
tight-loop code the cache hides this completely; for code with a large or scattered
footprint, misses are the single biggest placement cost on this chip.

\section{Putting hot code in SRAM}\index{Linker\_Section}

The escape hatch is to place code in IRAM, where it is never a cache miss because it
is not in flash at all. The runtime already does this for the things that must be
fast or must not fault during a flash stall: the exception vectors and the lowest
boot code sit in the \texttt{.iram1} sections the linker routes to SRAM. Application
code can join them with an aspect:

\begin{lstlisting}
procedure Hot_Path with Linker_Section => ".iram1";
\end{lstlisting}

IRAM is scarce (32\,KiB of it), so this is a scalpel, not a default: reserve it
for an interrupt handler or an inner loop that profiling shows is cache-bound. One
hard constraint bounds how much you can move: the 2nd-stage bootloader expects the
image to present exactly two flash-mapped segments (IROM and DROM), so the bulk of
the code must stay in flash; IRAM is for the hot slivers, not the program.

\section{PSRAM: large, slower, and DMA only on its terms}\index{PSRAM}\index{DMA}

External PSRAM gives you megabytes where SRAM gives you kilobytes, but it comes with
two performance caveats. It is reached through the cache and the octal SPI bus, so it
is slower than internal SRAM: fine for an interpreter's heap arena or a filesystem
block cache, not where you want a hot inner loop's working set. And DMA reaches it
only on strict terms. The GDMA controller sees PSRAM \emph{through the data
cache}, so a PSRAM buffer must be cache-line (32-byte) aligned, and ideally
whole-line sized. The GDMA driver then keeps it coherent, writing back and
invalidating the cache around each transfer, and the LCD driver streams its
framebuffers from PSRAM exactly this way. Internal SRAM needs none of that
ceremony and is faster besides. So small, hot scratch buffers stay internal:
the SD and filesystem drivers do exactly this, their scratch buffers internal
even when the rest of the data structure \textemdash{} or the task's own stack
(\texttt{ENV\_STACK\_PSRAM}, Chapter~\ref{ch:static-mem}) \textemdash{} is not.

\section{Measuring, not guessing}\index{CCOUNT}\index{cycle counter}

The CPU runs at 240\,MHz (the PLL that \texttt{start.S} selects out of the 20\,MHz
reset clock). At that rate, guessing about microseconds is hopeless; measure with the
cycle counter. The Xtensa \texttt{CCOUNT} register is a free-running 32-bit counter
at the core clock, readable in a single instruction. The SIMD benchmark
(Chapter~\ref{ch:asm}) uses exactly this to time its kernels:

\begin{lstlisting}
function Cycles return Unsigned_32 is
   V : Unsigned_32;
begin
   Asm ("rsr.ccount %0",
        Outputs => Unsigned_32'Asm_Output ("=r", V), Volatile => True);
   return V;
end Cycles;
\end{lstlisting}

Bracket a section between two reads, subtract (unsigned, so wrap takes care of
itself), and divide by 240 for microseconds. This is how every performance claim in
this book was obtained, including the discovery that an unplaced hot loop was
cache-bound, or that a SIMD kernel really was an order of magnitude faster than the
scalar one.

\section{Per-task CPU time}\index{Execution\_Time}\index{run statistics}\index{CPU usage}

\texttt{CCOUNT} times a code section; it cannot tell you where a whole
\emph{system}'s cycles go.  For that the runtime keeps per-thread execution
time: every context switch charges the elapsed slice to the outgoing thread
(\texttt{System.BB.Execution\_Time}), and interrupt time is accounted
separately.  \texttt{Thread\_Clock} returns any thread's accumulated time;
\texttt{Global\_Interrupt\_Clock} the handlers' total.  Both tick in the same
unit as the wall clock, so a task's share of an interval is a plain ratio.

Arming it costs one \texttt{with}.  The package's elaboration installs the
\texttt{Scheduling\_Event} hook, so a build whose closure never names it pays
nothing \textemdash{} the accounting simply does not exist.  The engine-monitor
example in the companion collection turns this into a run-statistics harness
(\texttt{ENGMON\_DEBUG=yes}): every five seconds, one line per registered task
with its CPU share, core affinity and stack high-water, and a closing line for
interrupt time:

\begin{lstlisting}[style=sh]
task env        c0  cpu 30.3%  stack 12288/hw 5036 (40%)
task decode     c1  cpu 1.3%   stack 20480/hw 11036 (53%)
task idle1      c1  cpu 98.3%  stack 8192/hw 384 (4%)
task interrupts     cpu 56.9%
\end{lstlisting}

The columns audit themselves: on the dual-core S3 the CPU column plus the
interrupt line must sum to about 200\% \textemdash{} both cores fully
accounted \textemdash{} and a pinned idle-counter task's share must match the
per-core idle measurement taken independently.  When a table like this
balances, you can trust what it says about the tasks in between.

\section{Optimisation levels}\index{optimisation}

The build is not uniform about \texttt{-O}, on purpose. The runtime, the HAL and the
libraries compile at \texttt{-O2} -- the scheduler and interrupt paths want the speed,
and so does anything you are benchmarking. The bare boot glue compiles at
\texttt{-Os} because bootloader size matters more than its speed and it runs once.
The host verification builds use \texttt{-O0 -g} so the debugger sees everything.
When you quote a number, quote the level with it: a \texttt{-O0} measurement of a
hot loop is not telling you what the shipped \texttt{-O2} build will do, and the gap
is often large.

\section{A placement story: the full-profile trampoline}

One placement problem is subtle enough to be worth seeing in full because it ties
the map back to the language. GCC implements a nested subprogram that captures its
enclosing scope with a \emph{trampoline}: a few bytes of code it writes onto the
stack. But the stack is in DRAM, a data-bus view with no execute permission, so
calling that trampoline faults. The \texttt{light-tasking} and \texttt{embedded}
profiles forbid the construct outright (\texttt{No\_Implicit\_Dynamic\_Code}); the
\texttt{full} profile instead, at task creation, detects a trampoline address in the
DRAM stack range and re-points it to the \emph{instruction-bus alias} of the same
physical SRAM: the identical bytes seen through an executable window
(Chapter~\ref{ch:runtime-build}). It is the memory map's defining feature (SRAM
visible at two addresses, one data and one instruction) turned from a hazard into
the fix.
